\documentclass[11pt,letterpaper]{article}
\usepackage[margin=0.85in]{geometry}
\usepackage{fontspec}
\setmainfont{TeX Gyre Pagella}
\usepackage{parskip}
\usepackage[hidelinks]{hyperref}
\pagestyle{empty}
\begin{document}
{\large\bfseries NP Labs at TRECVID 2026:\\Experiments with Off-the-Shelf Video QA Models\par}
\vspace{8pt}
\textbf{Task:} Video Question Answering (VQA)\\
\textbf{Group:} NP Labs\\
\textbf{Speaker:} Tanwir Zaman\\
\textbf{Affiliation:} Independent researcher\\
\textbf{Email:} \href{mailto:zaman.tanwir@gmail.com}{zaman.tanwir@gmail.com}\\
\textbf{Presentation mode:} In person
\vspace{8pt}

\textbf{Abstract}
\raggedright

Our TRECVID 2026 VQA participation explored how to build video question-answering submissions using open-weight Qwen models and paid Gemini inference, without task-specific training.

\textbf{Answer generation.} Across four runs, we tried sampled frames, speech transcripts, audio-video inputs, and combinations of answers from different models.

\textbf{Multiple choice.} Four runs explored a practical choice: reuse an answer generated earlier, or return to the video with the answer options available.

The talk will present these approaches and their official results, with a focus on where the submitted runs succeeded and failed. We will discuss question types that remained difficult, including counting, and use selected examples to compare answers from open-weight and paid models. These examples will illustrate differences between the approaches and motivate further experiments.
\end{document}
